\section{Randomised SVD Preprocessing}\label{sec:preprocessing}

The Woodbury solver requires the top-$k$ eigenpairs of the sample correlation
$C = D^{-1/2}\hat\Sigma\,D^{-1/2}$, equivalently of the standardised returns
$X_s = X D^{-1/2}$. Two paths are available.

\paragraph{Dense path}
Assemble $C$ explicitly ($O(n^2 T)$ FLOPs, $O(n^2)$ storage), then compute a full
eigendecomposition ($O(n^3)$). Total cost $O(n^2 T + n^3)$; $0.011$\,s at $n = 494$, $T = 1213$ on the benchmark
hardware.

\paragraph{Randomised SVD path}
The right singular vectors of $X_s/\sqrt{T}$ are the eigenvectors of $C$. A
randomised truncated SVD~\cite{halko2011} computes the top-$k$ eigenpairs without
forming $X_s^\top X_s$. The trace-preserving bulk mean is recovered without any
dense decomposition from
$\bar\mu = \bigl(n - \sum_{j\le k}\hat\lambda_j\bigr)/(n-k)$, using
$\mathrm{tr}(C) = n$.

\begin{algorithm}[ht]
\caption{Randomised truncated SVD for eigenpairs of the correlation $C$}\label{alg:rsvd}
\begin{algorithmic}[1]
\Require Standardised returns $X_s = X D^{-1/2} \in \mathbb{R}^{T \times n}$; target rank $k$; oversampling $p \geq 5$
\Ensure  Signal eigenvectors $U_k \in \mathbb{R}^{n \times k}$,
         eigenvalues $\Lambda_k \in \mathbb{R}^k$, bulk mean $\bar\mu$
\State $\Omega \leftarrow \mathcal{N}(0, I_{n \times (k+p)})$ \Comment{random test matrix}
\State $Y \leftarrow X_s^\top(X_s \Omega)$ \Comment{$n \times (k+p)$; two matrix-vector products with $X_s$, $O(nT(k+p))$}
\State $Q \leftarrow \mathrm{orth}(Y)$ \Comment{orthonormal basis, $O(n(k+p)^2)$}
\State $B \leftarrow (X_s Q)^\top(X_s Q) / T$ \Comment{$(k+p) \times (k+p)$ projected correlation; no $X_s^\top X_s$ formed}
\State $[\hat U, \hat\Lambda, \sim] \leftarrow \mathrm{eigh}(B)$; sort descending
\State $U_k \leftarrow Q \hat U[:, 1{:}k]$; $\Lambda_k \leftarrow \hat\Lambda[1{:}k]$
\State $\bar\mu \leftarrow \bigl(n - \textstyle\sum_{j\le k}\hat\Lambda_j\bigr)/(n-k)$; identify $k$ via MP threshold
\State \Return $(U_k, \Lambda_k, \bar\mu)$
\end{algorithmic}
\end{algorithm}

\begin{proposition}[Randomised SVD cost and accuracy]\label{prop:rsvd}
Algorithm~\ref{alg:rsvd} costs $O(nT(k+p))$ FLOPs and $O(n(k+p))$ storage; no
$n \times n$ matrix is formed. Let $P = QQ^\top$ denote the orthogonal
projector onto the computed basis. Since the sketch $Y = T\,\hat\Sigma\,\Omega$
applies the range finder of~\cite{halko2011} to the symmetric matrix
$\hat\Sigma$ with Gaussian test matrix $\Omega$ and oversampling $p \geq 2$,
Theorem~10.6 of~\cite{halko2011} gives, in expectation,
\[
  \mathbb{E}\,\bigl\lVert (I - P)\hat\Sigma \bigr\rVert_2
  \;\leq\;
  \Bigl(1 + \sqrt{\tfrac{k}{p-1}}\Bigr)\,\lambda_{k+1}
  \;+\; \frac{e\sqrt{k+p}}{p}\Bigl(\textstyle\sum_{j>k}\lambda_j^2\Bigr)^{1/2}
  \;=:\; \epsilon_{k,p}\,,
\]
where $\lambda_j$ are the eigenvalues of $\hat\Sigma$. The returned eigenpairs
are the exact eigenpairs of the compression $P\hat\Sigma P$, and by Weyl's
inequality the eigenvalue errors satisfy
$|\lambda_j - \hat\lambda_j| \leq \lVert\hat\Sigma - P\hat\Sigma P\rVert_2
 \leq 2\,\lVert (I-P)\hat\Sigma \rVert_2$ for all $j \leq k$.
\end{proposition}

\begin{proof}
The cost and storage counts follow from the line-by-line annotations of
Algorithm~\ref{alg:rsvd}. The error bound is Theorem~10.6
of~\cite{halko2011} applied to $A = \hat\Sigma$ (the deterministic scaling $T$
of the sketch does not change its range). For the compression bound, write
$\hat\Sigma - P\hat\Sigma P = (I-P)\hat\Sigma + P\hat\Sigma(I-P)$ and note that
both terms have spectral norm at most $\lVert(I-P)\hat\Sigma\rVert_2$ by
symmetry; Weyl's inequality~\cite{hornjohnson2013} then bounds each eigenvalue
perturbation by the norm of the difference.
\end{proof}

The bound $\epsilon_{k,p}$ is governed by the correlation eigenvalues at and below
the MP edge $\lambda_{k+1} \approx \hat\lambda_+ \approx 2.28$ on S\&P~500 data,
against a leading signal eigenvalue $\lambda_1 \approx 313\,\bar\mu$ that is two
orders of magnitude larger. The bound is conservative (the tail-sum term
accumulates all $n - k$ bulk eigenvalues) and the measured accuracy is far better.

\paragraph{Memory passes}
Algorithm~\ref{alg:rsvd} involves four matrix-vector products with $X$ in total:
two for the sketch $Y = X^\top(X\Omega)$ (line~2) and two for the projection
$B = Q^\top X^\top XQ$ (line~4, computed as $(XQ)^\top(XQ)/T$). Reading $X$
from memory is required at two distinct stages: once for the sketch (line~2)
and once for the projection (line~4). For very large $n$ and $T$ where $X$
does not fit in memory, a single-pass variant that combines both stages is
available at a modest increase in approximation error; see~\cite{halko2011}
for details. At benchmark size ($n=494$, $T=1213$, approximately $4.6$\,MB
for double-precision $X$) both stages fit comfortably in L3 cache.

\paragraph{Benchmark}
On S\&P~500 ($n=494$, $T=1213$, $k=17$, $p=10$) the randomised SVD is $1.9\times$
faster than the dense path ($0.006$\,s vs $0.011$\,s), needs $29\times$ less
storage, and aligns the top-$k$ signal subspace to within a Frobenius error
$\norm{UU^\top - \hat{U}\hat{U}^\top}_F \approx 2.5 \times 10^{-2}$ (about $0.6\%$
per eigenvector column), an order of magnitude below the $k \pm 1$ threshold
uncertainty (Section~\ref{sec:ksensitivity}) that any user of the MP rule already
accepts. The advantage grows with $n$: at $n=3000$ the dense eigendecomposition
takes $1.12$\,s against $0.056$\,s for rSVD, a $20\times$ saving
(Figure~\ref{fig:scaling_full}).

\paragraph{Large-\texorpdfstring{$n$}{n} recommendation}
For $n \leq 500$ either path is practical; the dense path is simpler to implement
and both finish in milliseconds. For $n > 1000$ the randomised SVD is strongly
preferred: the dense path's $O(n^2)$ storage requirement alone (8\,MB at $n=1000$,
72\,MB at $n=3000$) begins to stress L3 cache, while the randomised path's
$O(n(k+p))$ requirement (about $3$\,MB at $n=3000$, $k=125$) stays within
typical L3.
